\section*{Chapitre \ref{ch:b1:chemical-biogenic-sediments} --- Les sédiments carbonatés, siliceux et évaporitiques}

\begin{solution}{exo:b1:chemical-biogenic-sediments:1}
Chlorure $0.5503\times35.165 = \qty{19.35}{g}$~; sodium $0.3066\times35.165 =
\qty{10.78}{g}$~; sulfate \qty{2.711}{g}~; magnésium \qty{1.284}{g}. Le sodium et le
chlorure font ensemble $0.5503 + 0.3066 = 0.857$, soit \qty{85.7}{\%} du
sel.
\end{solution}

\begin{solution}{exo:b1:chemical-biogenic-sediments:2}
La fraction restante est $1/\mathrm{CF}$~: carbonate de calcium $1/1.8 = 0.556$ restant,
\qty{44.4}{\%} évaporés~; gypse $0.263$ restant, \qty{73.7}{\%} évaporés~;
halite $0.094$ restant, \qty{90.6}{\%} évaporés~; sulfate de magnésium $0.014$ restant,
\qty{98.6}{\%} évaporés.
\end{solution}

\begin{solution}{exo:b1:chemical-biogenic-sediments:3}
(a)~$\rho_s(1-\phi) = G/a = 8.4/0.0042 = \qty{2000}{kg.m^{-3}}$. (b)~Avec
l'aragonite, $1 - 2000/2930 = 0.32$, soit une porosité de \qty{32}{\%}~; avec la calcite,
$1 - 2000/2710 = 0.26$. (c)~$a = 3.7/2000 = \qty{1.85e-3}{m}$, soit \qty{1.85}{mm} par
an.
\end{solution}

\begin{solution}{exo:b1:chemical-biogenic-sediments:4}
Jusqu'à \qty{4000}{m} la teneur ne diminue que lentement (88 à 80)~; en dessous de
\qty{4000}{m} elle chute brutalement~: la \omterm{def:b1:chemical-biogenic-sediments:ccd}{lysocline} se situe vers \qty{4000}{m}. La
teneur passe sous \qty{10}{\%} entre \qty{4700}{m} (12) et \qty{4800}{m}
(6)~: $4700 + 100\times(12-10)/(12-6) = \qty{4733}{m}$, soit une PCC à environ
\qty{4730}{m}, quelque \qty{700}{m} sous la \omterm{def:b1:chemical-biogenic-sediments:ccd}{lysocline}.
\end{solution}

\begin{solution}{exo:b1:chemical-biogenic-sediments:5}
(a)~\ce{Ca^2+ + 2 HCO3- -> CaCO3 + CO2 + H2O}~: une mole de \ce{CO2} par mole de
\ce{CaCO3}, donc $44.01/100.09 = \qty{0.440}{kg}$ de \ce{CO2} par kilogramme de
\ce{CaCO3}. (b)~$3.7\times0.440 = \qty{1.63}{kg}$ de \ce{CO2} par mètre carré et
par an. (c)~Sur les deux ions hydrogénocarbonate utilisés, un seul carbone entre dans
le \omterm{def:b1:rock-forming-minerals:mineral}{minéral}~; l'autre part sous forme de dioxyde de carbone. La calcification retire de l'alcalinité
à l'eau et abaisse sa capacité à retenir du carbone dissous~; la part du
\ce{CO2} qui atteint l'air dépend du pouvoir tampon de l'eau de mer, affaire des
équilibres des carbonates des livres de chimie.
\end{solution}

\begin{solution}{exo:b1:chemical-biogenic-sediments:6}
(a)~$t_c = (2000/350)^2 = \qty{32.7}{\Ma}$~; à \qty{25}{km} par million d'années,
$25\times32.7 = \qty{816}{km}$ de l'axe. (b)~$t_c = (2600/350)^2 =
\qty{55.2}{\Ma}$, \qty{1380}{km}. (c)~Un fond de \qty{60}{\Ma} a reçu de l'argile
rouge pendant $60 - 32.7 = \qty{27.3}{\Ma}$ dans le premier cas, et pendant seulement
$60 - 55.2 = \qty{4.8}{\Ma}$ dans le second.
\end{solution}

\begin{solution}{exo:b1:chemical-biogenic-sediments:7}
(a)~Oui~: $a = 4.2 > r = 3$, et le \omterm{def:b1:chemical-biogenic-sediments:reef}{récif} gagne \qty{1.2}{mm} par an sur la mer
jusqu'à ce que son sommet atteigne la surface. (b)~$D$ croît de $10 - 4.2 = \qty{5.8}{mm}$ par
an~; de 2 à \qty{50}{m} il faut $48/0.0058 = \num{8300}$ ans. (c)~Parce que
sa vitesse de croissance diminue à mesure que l'eau s'approfondit~: moins de lumière parvient aux algues qui
vivent dans les coraux et les nourrissent, si bien que $a$ décroît et que le \omterm{def:b1:chemical-biogenic-sediments:reef}{récif} prend du retard
de plus en plus vite.
\end{solution}

\begin{solution}{exo:b1:chemical-biogenic-sediments:8}
(a)~La réaction est \ce{CaSO4.2H2O -> CaSO4 + 2 H2O}. (b)~Les volumes molaires sont $172.17/2.30 =
\qty{74.9}{cm^3}$ et $136.14/2.97 = \qty{45.8}{cm^3}$, donc le solide perd
$1 - 45.8/74.9 = \qty{38.8}{\%}$ de son volume. (c)~$2\times18.02/172.17 = 0.209$~:
\qty{20.9}{\%} de la masse part sous forme d'eau. Libérée en profondeur, cette eau
élève la pression dans les pores des roches environnantes~; sous un couvercle étanche elle peut affaiblir
les roches et aider les failles à glisser, et elle peut dissoudre et emporter d'autres
sels.
\end{solution}

\begin{solution}{exo:b1:chemical-biogenic-sediments:9}
(a)~$9.2/255 = \qty{3.6}{\%}$ préservés, \qty{96.4}{\%} dissous. (b)~Là où
la production est très forte, quelques pour cent représentent encore un grand flux, et là où
peu d'argile ou de carbonate la dilue~: sous les eaux de remontée de l'océan Austral,
du Pacifique Nord et du Pacifique équatorial, et sous la PCC.
(c)~Les frustules minces et délicats se dissolvent d'abord et les robustes survivent~: le
sédiment garde une sélection des espèces les plus résistantes, non la communauté qui
vivait dans l'eau.
\end{solution}

\begin{solution}{exo:b1:chemical-biogenic-sediments:10}
(a)~$EA = 1.5\times\qty{1e10}{m^2} = \qty{1.5e10}{m^3}$ par an. (b)~Avec
$Q_{\mathrm{in}} = \qty{1.8e10}{m^3}$, $Q_{\mathrm{out}} = \qty{0.3e10}{m^3}$ et
$k = 6$~: le seuil du gypse (3.8) est dépassé mais pas celui de la halite (10.6), donc
le carbonate de calcium et le gypse précipitent. Avec $Q_{\mathrm{in}} =
\qty{1.575e10}{m^3}$, $Q_{\mathrm{out}} = \qty{0.075e10}{m^3}$ et $k = 21$, pour
un ion qui reste en solution tel que le magnésium~: la saumure dépasse le seuil de la
halite, donc la halite précipite aussi, mais pas les sels de magnésium (70).
(c)~La halite exige $Q_{\mathrm{in}}/(Q_{\mathrm{in}} - EA) \ge 10.6$, c'est-à-dire
$Q_{\mathrm{in}} \le 10.6\,EA/9.6 = 1.104\,EA = \qty{1.66e10}{m^3}$ par an~:
l'apport ne peut dépasser l'évaporation que d'un dixième environ au plus.
\end{solution}

\begin{solution}{exo:b1:chemical-biogenic-sediments:11}
(a)~L'eau de mer contient \qty{67.3}{mg} de bromure par kilogramme~; à 10.6~fois, la
saumure en contient $67.3\times10.6 = \qty{713}{mg.kg^{-1}}$ et la première halite
$0.032\times713 = \qty{22.8}{mg.kg^{-1}}$. (b)~À 30~fois, $67.3\times30 =
\qty{2019}{mg.kg^{-1}}$ dans la saumure et \qty{64.6}{mg.kg^{-1}} dans la halite.
(c)~La halite ne prend que 3.2~centièmes de la concentration de la saumure, de sorte que
presque tout le bromure reste en solution et continue de se concentrer à mesure que l'eau
s'évapore. (d)~La hausse enregistre la concentration de la saumure~; la chute brusque
enregistre sa dilution, par un nouvel apport d'eau de mer dans le bassin, ou d'eau
douce, avant que la halite continue de croître.
\end{solution}

\begin{solution}{exo:b1:chemical-biogenic-sediments:12}
(a)~$t_c = \qty{32.65}{\Ma}$, donc $0.02\times\num{32650} = \qty{653}{m}$ de
vase (\num{32650} intervalles de mille ans à \qty{2}{cm}) et
$0.001\times\num{37350} = \qty{37}{m}$ d'argile rouge. (b)~$t_c = (2500/350)^2 =
\qty{51.0}{\Ma}$~: \qty{1020}{m} de vase et \qty{19}{m} d'argile rouge. Une PCC plus profonde
donne une vase bien plus épaisse, et l'argile rouge du fond profond de l'Atlantique est plus mince
que celle du Pacifique du même âge. (c)~La PCC a pu être plus profonde pendant un
temps, si bien que le site est repassé au-dessus d'elle~; ou la plaque a porté le site sous
une ceinture de forte productivité, comme l'équateur, où davantage de carbonate pleuvait~;
ou la couche a été apportée par un glissement ou un \omterm{def:b1:clastic-sediments:turbidite}{courant de turbidité} depuis une partie
plus haute du fond marin.
\end{solution}

\begin{solution}{pb:b1:chemical-biogenic-sediments:1}
\textbf{1.} $1030\times1000 = \qty{1.03e6}{kg}$ d'eau de mer, contenant
$1.03\times10^6\times0.035165 = \qty{36200}{kg}$ de sel.

\textbf{2.} Gypse $1.03\times10^6\times1.581\times10^{-3} = \qty{1628}{kg}$~;
halite $1.03\times10^6\times0.02741 = \qty{28230}{kg}$.

\textbf{3.} Carbonate de calcium $113.3/2710 = \qty{0.042}{m}$~; gypse $1628/2300 =
\qty{0.71}{m}$~; halite $28230/2170 = \qty{13.0}{m}$.

\textbf{4.} L'anhydrite a la masse $1628\times136.14/172.17 =
\qty{1288}{kg}$ et l'épaisseur $1288/2970 = \qty{0.43}{m}$.

\textbf{5.} À sec, le sel fait $0.110 + 1.250 + 27.41 + 6.34 = \qty{35.11}{g}$ par
kilogramme, dont la halite fait $27.41/35.11 = \qty{78}{\%}$~: le sodium et le chlorure
font l'essentiel du sel dissous, et la halite en prend presque la totalité.

\textbf{6.} De bas en haut~: carbonate de calcium, gypse, halite, sels
d'amertume. Chaque \omterm{def:b1:rock-forming-minerals:mineral}{minéral} commence à précipiter quand la saumure atteint sa
saturation, le moins soluble d'abord, et se dépose sur ce qui l'a précédé.

\textbf{7.} La profondeur de saumure est $1000/\mathrm{CF}$~: \qty{556}{m} pour le carbonate de
calcium, \qty{263}{m} pour le gypse, \qty{94}{m} pour la halite.

\textbf{8.} $1 - 1/10.6 = \qty{90.6}{\%}$.

\textbf{9.} $1000/70 = \qty{14}{m}$ pour le sulfate de magnésium, $1000/90 = \qty{11}{m}$
pour les sels de potassium.

\textbf{10.} Environ \qty{1.10}{g.cm^{-3}} quand le carbonate de calcium se forme, 1.13 avec
le gypse, 1.22 avec la halite et plus de 1.30 avec les sels d'amertume, contre 1.03 pour
l'eau de mer. L'eau de mer qui afflue flotte sur la saumure dense, qui reste au
fond~; dans un bassin barré la saumure peut retourner à l'océan sous
l'apport.

\textbf{11.} Il n'a que \qty{4}{cm} d'épaisseur, contre \qty{13}{m} de halite~; et il
se forme sur les marges, où l'eau qui afflue se concentre d'abord, plus qu'au
centre, si bien qu'il se mêle à la base du gypse ou est remplacé par lui.

\textbf{12.} La halite ne prend qu'une petite part du bromure de la saumure
(exercice~\ref{exo:b1:chemical-biogenic-sediments:11}), donc la saumure devient
plus riche en bromure à mesure que la halite croît, et chaque nouvelle couche de halite en capte
davantage.

\textbf{13.} $1500/13.01 = 115.3$~: 116 remplissages.

\textbf{14.} $116\times0.434 = \qty{50}{m}$ d'anhydrite, non en un seul banc mais en
116 couches, chacune un peu moins d'un demi-mètre d'épaisseur, à la base de la halite de
son remplissage~: une succession cyclique d'anhydrite et de halite.

\textbf{15.} Par remplissage, $1.03\times10^6\times6.34\times10^{-3}/2000 =
\qty{3.3}{m}$, donc \qty{379}{m} pour 116 remplissages. Ce sont les sels les plus
solubles~: l'apport suivant d'eau de mer les redissout, et un bassin dont la saumure
ressort avant la fin n'y parvient jamais.

\textbf{16.} À chaque remplissage la saumure est diluée~: la teneur en bromure chute
brutalement à la base de la nouvelle halite puis remonte vers le haut, un profil en dents de
scie avec une dent par remplissage.

\textbf{17.} $2000/13.01 = 154$ remplissages.

\textbf{18.} $5000\times365.25\times86400 = \qty{1.578e11}{m^3}$ par an, répartis
sur $\qty{1.0e11}{m^2}$~: $E = \qty{1.58}{m}$ par an.

\textbf{19.} $1.578\times10^{11}\times1030\times0.035165 = \qty{5.72e12}{kg}$ de
sel par an.

\textbf{20.} $1.578\times10^{11}\times1030\times0.02741/2170 =
\qty{2.05e9}{m^3}$ de halite par an, soit \qty{0.0205}{m}, environ \qty{2}{cm}
par an sur le fond.

\textbf{21.} $1500/0.0205 = \num{73000}$ ans.

\textbf{22.} Le niveau baisse de \qty{1.58}{m} par an~; \qty{1000}{m} d'eau ont
disparu en $1000/1.578 = 634$ ans.

\textbf{23.} La halite apparaît quand \qty{90.6}{\%} de l'eau est partie, au bout de
$905.7/1.578 = 574$ ans.

\textbf{24.} Les \qty{1.5}{km} de halite ont pris le sel de 116 remplissages du
bassin (115.3, à arrondir au-dessus), qu'ils soient venus l'un après l'autre à mesure que le
bassin s'asséchait et se remplissait, ou en continu par-dessus le seuil, comme à la
question~21, au rythme d'un bassin d'eau de mer tous les
$1000/1.578 = 634$ ans.
\end{solution}
